\documentclass[compress, 8pt, aspectratio=169]{beamer}

\usepackage{presentationtemplate}
\usepackage[askip=3mm, bskip=3mm]{terminal}
\usepackage[linenosfontsize=\tiny, askip=3mm, bskip=3mm]{mylisting}
\tcbuselibrary{hooks}
\usepackage{csquotes}

\title{Стандартная библиотека: строки}

\begin{document}

  \frame[plain]{\titlepage}

  \begin{frame}[fragile]

    \frametitle{Строки: \texttt{char[]} vs \texttt{std::string}}

    \begin{itemize}

      \item \texttt{char[]}: изменение посимвольно;
        операторов сравнения и конкатенации нет; память \textendash \space вручную.

      \item \texttt{std::string}: значение; память управляет объект;
        строку можно переприсвоить целиком.

    \end{itemize}

    \myinputlisting[minted language=cpp, minted options app={fontsize=\footnotesize}]
        {Presentations/18-strings/}
        {why.cpp}

    \begin{terminalwindow}
!\shellcommand{g++ why.cpp -o why \&\& ./why}!
Eda
Grace
    \end{terminalwindow}

  \end{frame}

  \begin{frame}[fragile]

    \frametitle{Создание строк}

    \begin{itemize}

      \item Инициализация: из литерала, из \texttt{(count, char)},
        копирование, перемещение.

      \item \texttt{(count, char)} \textendash \space повтор одного символа.

    \end{itemize}

    \myinputlisting[minted language=cpp]
        {Presentations/18-strings/}
        {create.cpp}

    \begin{terminalwindow}
!\shellcommand{g++ create.cpp -o create \&\& ./create}!
C++ ***** C++
    \end{terminalwindow}

  \end{frame}

  \begin{frame}[fragile]

    \frametitle{Конкатенация и сравнение}

    \begin{itemize}

      \item Конкатенация: \texttt{+}, \texttt{+=} (также \texttt{append}).

      \item Сравнение: \texttt{<}, \texttt{==}, \texttt{<=>}
        \textendash \space лексикографическое, побайтово.

    \end{itemize}

    \myinputlisting[minted language=cpp]
        {Presentations/18-strings/}
        {concat-compare.cpp}

    \begin{terminalwindow}
!\shellcommand{g++ concat-compare.cpp -o concat-compare \&\& ./concat-compare}!
true true
    \end{terminalwindow}

  \end{frame}

  \begin{frame}[fragile]

    \frametitle{Доступ к символам}

    \begin{itemize}

      \item \texttt{[]} \textendash \space без проверки границ (UB при выходе за них);
        \texttt{at} \textendash \space бросает \texttt{std::out\_of\_range}.

      \item \texttt{front}/\texttt{back};
        \texttt{c\_str()}/\texttt{data()} \textendash \space строка как C-строка
        (с завершающим нулём).

    \end{itemize}

    \myinputlisting[minted language=cpp]
        {Presentations/18-strings/}
        {access.cpp}

    \begin{terminalwindow}
!\shellcommand{g++ access.cpp -o access \&\& ./access}!
C
+
+
    \end{terminalwindow}

  \end{frame}

  \begin{frame}[fragile]

    \frametitle{Поиск в строке}

    \begin{itemize}

      \item \texttt{find}/\texttt{rfind} \textendash \space позиция вхождения
        или \texttt{std::string::npos}, если не найдено.

      \item \texttt{starts\_with}/\texttt{ends\_with} \textendash \space C++20;
        \texttt{contains} \textendash \space C++23.

    \end{itemize}

    \myinputlisting[minted language=cpp, minted options app={fontsize=\footnotesize}]
        {Presentations/18-strings/}
        {find.cpp}

    \begin{terminalwindow}
!\shellcommand{g++ find.cpp -o find \&\& ./find}!
16
17
true true true
    \end{terminalwindow}

  \end{frame}

  \begin{frame}[fragile]

    \frametitle{Подстроки: \texttt{substr}, \texttt{replace}, \texttt{erase}}

    \begin{itemize}

      \item \texttt{substr(pos, len)} \textendash \space копия подстроки
        (без \texttt{len} \textendash \space до конца).

      \item \texttt{replace(pos, len, str)} \textendash \space замена участка;
        \texttt{erase(pos, len)} \textendash \space удаление участка.

    \end{itemize}

    \myinputlisting[minted language=cpp]
        {Presentations/18-strings/}
        {modify.cpp}

    \begin{terminalwindow}
!\shellcommand{g++ modify.cpp -o modify \&\& ./modify}!
world
bye, world
bye,
    \end{terminalwindow}

  \end{frame}

  \begin{frame}[fragile]

    \frametitle{Преобразование с числами}

    \begin{itemize}

      \item В строку: \texttt{std::to\_string} (числовые типы).

      \item Из строки: \texttt{stoi}, \texttt{stoul}, \texttt{stof},
        \texttt{stod}, ...; третий аргумент \textendash \space основание.

      \item Ошибка парсинга \textendash \space исключение
        \texttt{invalid\_argument}/\texttt{out\_of\_range}.

    \end{itemize}

    \myinputlisting[minted language=cpp]
        {Presentations/18-strings/}
        {convert.cpp}

    \begin{terminalwindow}
!\shellcommand{g++ convert.cpp -o convert \&\& ./convert}!
42
3.14
255
100
    \end{terminalwindow}

  \end{frame}

  \begin{frame}[fragile]

    \frametitle{Small String Optimization (SSO)}

    \begin{itemize}

      \item Короткие строки хранятся в самом объекте (на стеке),
        без выделения в куче.

      \item Порог \textendash \space свойство реализации: libstdc++ (x86-64)
        \textendash \space 15 символов, \texttt{sizeof(std::string)} = 32.

      \item Проверяем по адресу \texttt{data()}: короткий
        \textendash \space внутри объекта, длинный \textendash \space в куче.

    \end{itemize}

    \myinputlisting[minted language=cpp, minted options app={fontsize=\footnotesize}]
        {Presentations/18-strings/}
        {sso.cpp}

    \begin{terminalwindow}
!\shellcommand{g++ sso.cpp -o sso \&\& ./sso}!
32
2: 0x7ffd4854fd00
42: 0x5c2d7fb61360
    \end{terminalwindow}

  \end{frame}

  \begin{frame}[fragile]

    \frametitle{Unicode: строка \textendash \space это байты}

    \begin{itemize}

      \item \texttt{std::string} хранит байты; по соглашению \textendash \space UTF-8.

      \item \texttt{size()}/\texttt{length()} \textendash \space число байт,
        а не символов.

      \item \texttt{s[i]} \textendash \space по байтам: многобайтовый символ
        \enquote{разрывается}.

    \end{itemize}

    \myinputlisting[minted language=cpp]
        {Presentations/18-strings/}
        {utf8-bytes.cpp}

    \begin{terminalwindow}
!\shellcommand{g++ utf8-bytes.cpp -o utf8-bytes \&\& ./utf8-bytes}!
bytes: 6
f0 9f 9a 80 68 69
    \end{terminalwindow}

  \end{frame}

  \begin{frame}[fragile]

    \frametitle{Обход по символам: \texttt{std::u8string\_view}}

    \begin{itemize}

      \item \texttt{std::u8string\_view} (C++20) \textendash \space вид
        на UTF-8-байты; итерация идёт по байтам \texttt{char8\_t},
        не по кодовым точкам.

      \item Конструктора из \texttt{std::string} нет: делаем вид
        на байты строки \textendash \space \enquote{строка \textendash \space просто байты}.

      \item Число кодовых точек \textendash \space число непродолжающих
        байт (не начинающихся с \texttt{0b10}).

    \end{itemize}

    \myinputlisting[minted language=cpp, minted options app={fontsize=\footnotesize}]
        {Presentations/18-strings/}
        {u8string-view.cpp}

  \end{frame}

  \begin{frame}[fragile]

    \frametitle{Вывод программы с листинга на прошлом слайде}

    \begin{terminalwindow}
!\shellcommand{g++ u8string-view.cpp -o u8string-view \&\& ./u8string-view}!
bytes: 12 chars: 6
    \end{terminalwindow}

  \end{frame}

  \begin{frame}[fragile]

    \frametitle{Строковый литерал \texttt{""s} (C++14)}

    \begin{itemize}

      \item Литерал \texttt{"...s"} \textendash \space объект
        \texttt{std::string}; \texttt{"...sv"} (C++17)
        \textendash \space \texttt{std::string\_view}.

      \item Каждый литерал \textendash \space новое временное значение:
        \texttt{"a"s + "b"} \textendash \space два временных \texttt{std::string}.

      \item \texttt{using namespace std::literals::string\_literals}
        \textendash \space стандартный способ активации.

    \end{itemize}

    \myinputlisting[minted language=cpp]
        {Presentations/18-strings/}
        {s-literal.cpp}

    \begin{terminalwindow}
!\shellcommand{g++ s-literal.cpp -o s-literal \&\& ./s-literal}!
C++ 23
    \end{terminalwindow}

  \end{frame}

\end{document}
